% CS 263 Phase 0 feasibility report.
% Build with `make report` from the repository root. Every number comes from
% build/report/generated/, written by analysis/report_assets.py from the runs
% listed in analysis/report_runs.json; this file types no measured values.
\documentclass[11pt]{article}

% Reproducible PDF: no build timestamp or random trailer ID.
\pdfinfoomitdate=1
\pdftrailerid{}
\pdfsuppressptexinfo=-1

\usepackage[T1]{fontenc}
\usepackage[utf8]{inputenc}
\usepackage{lmodern}
\usepackage[margin=1in]{geometry}
\usepackage{microtype}
\usepackage{booktabs}
\usepackage{graphicx}
\usepackage{float}
\usepackage{xcolor}
\usepackage[font=small,labelfont=bf]{caption}
\usepackage[hidelinks]{hyperref}
\usepackage{enumitem}
\setlist{itemsep=2pt,topsep=4pt}

% Generated assets, relative to report/ (latexmk runs with -cd).
\providecommand{\GeneratedDir}{../build/report/generated}
\input{\GeneratedDir/numbers.tex}

\title{CPython Engines on gem5: Phase 0 Feasibility Report}
\author{CS 263 project}
\date{}

\begin{document}
\maketitle

\noindent CPython \CPythonTag{} runs in gem5 under all four engine configurations, and the JIT demonstrably executes in the simulated hot loop.
On a small smoke workload, the tail-call interpreter took \CycleSavingPctTailCall\% fewer cycles than the default interpreter and the JIT \CycleSavingPctJitOn\% fewer.
These are single feasibility runs, not findings: startup is included and a JIT probe sits inside the measured loop (Section~\ref{sec:limits}).

\section{Setup}

Every engine is built from one CPython commit with one compiler and runs on the same simulated x86-64 core.
gem5 runs natively in an arm64 container on Apple silicon and simulates x86-64.

\begin{table}[H]
\centering\small
\input{\GeneratedDir/tab-setup.tex}
\caption{Pinned inputs.}
\end{table}

\paragraph{Engine configurations.} Three builds give four configurations:
\begin{itemize}
  \item \emph{default}: specializing interpreter, computed-goto dispatch;
  \item \emph{tail-call}: configured with \texttt{-{}-with-tail-call-interp};
  \item \emph{JIT off} and \emph{JIT on}: one binary configured with \texttt{-{}-enable-experimental-jit=yes-off}, run with \texttt{PYTHON\_JIT=0} or \texttt{PYTHON\_JIT=1}.
\end{itemize}

\paragraph{Workload.} A Python Collatz loop over 1--2000 (\InnerIterations{} inner iterations) whose output line is checked against an independent C implementation.
The guest environment holds only two variables: \texttt{PYTHON\_JIT} and \texttt{PYTHONHASHSEED=0}.

\section{Feasibility by engine}

All four configurations pass every check: gem5 exited normally, the guest exited 0, the output matched, and all required counters were recorded.
The engine mode was confirmed from inside the guest, not inferred from the build (Table~\ref{tab:feas}).

\begin{table}[H]
\centering\small
\input{\GeneratedDir/tab-feasibility.tex}
\caption{Per-engine checks and mode evidence.}
\label{tab:feas}
\end{table}

\paragraph{The JIT works under syscall emulation.} gem5 ignores the \texttt{mprotect} call the JIT uses to make its code executable (\MprotectJit{} warnings in the JIT run, \MprotectOther{} in the default run).
Even so, \JitActiveJitOn{} of \InnerIterations{} hot-loop iterations (\JitActivePct\%) ran inside JIT code in gem5, the same count as under Docker's amd64 emulation outside gem5.

\section{Results}

\begin{figure}[H]
\centering
\includegraphics[width=\textwidth]{\GeneratedDir/fig-relative.pdf}
\caption{Tail-call saves cycles; the JIT saves instructions. Whole run relative to the default interpreter (default $=1.00$), O3 baseline, one run each.}
\label{fig:relative}
\end{figure}

The tail-call interpreter runs nearly the same instructions as default (\RelInstsTailCall) but needs fewer cycles (\RelCyclesTailCall).
That is consistent with its \BrCutPctTailCall\% fewer branch mispredictions (Table~\ref{tab:counters}); the mechanism is a proposal, not yet tested.
The JIT's gain comes mostly from running fewer instructions (\RelInstsJitOn), at lower IPC.
The JIT-capable binary with the JIT off runs \InstExtraPctJitOff\% more instructions than default for about the same cycles.

\begin{figure}[H]
\centering
\includegraphics[width=\textwidth]{\GeneratedDir/fig-startup.pdf}
\caption{Startup is about \StartupMeanM M instructions in every engine, \StartupShareMin--\StartupShareMax\% of each run. Startup measured with input 1 on gem5's atomic CPU; workload $=$ total minus startup.}
\label{fig:startup}
\end{figure}

The JIT shrinks the workload from \WorkMDefault M to \WorkMJitOn M instructions (\JitWorkSavingPct\% fewer) while startup is unchanged.
The subtraction relies on O3 and the atomic CPU retiring the same instruction stream: \InstStreamNote.
Startup cycles on O3 were not measured, so cycles cannot be split the same way.

\begin{table}[H]
\centering\small
\input{\GeneratedDir/tab-counters.tex}
\caption{Engine counters on the O3 baseline.}
\label{tab:counters}
\end{table}

\section{C smoke test: branch predictor sweep}

Changing one hardware parameter works as intended.
On a static C Collatz loop, resizing the predictor's global table moved cycles by at most \CSweepMaxDeltaPct\%.
The loop's key branch tests whether a number is odd, which is close to random, so \CSweepMPKIMin--\CSweepMPKIMax{} mispredictions per 1k instructions remain at every size.

\begin{table}[H]
\centering\small
\input{\GeneratedDir/tab-sweep.tex}
\caption{One-factor sweep of \texttt{globalPredictorSize}. Every run executed exactly \CSweepInsts{} instructions.}
\end{table}

Repeating the baseline reproduced an earlier run's simulated time to the tick (\CBaselineTicks{} ticks), and halving the L1 instruction cache gave the same tick count.

\section{Limits}
\label{sec:limits}

The engine ratios above show the pipeline works; they do not yet measure the engines.
\begin{itemize}
  \item \textbf{The probe is in the measured loop.} Every engine calls \texttt{sys.\_jit.is\_active()} on each hot-loop iteration. That adds about \ProbeOverheadM M instructions to the default run (\NoProbeInstsM M without it, \ProbeInstsM M with it), so the ratios partly measure that builtin call.
  \item \textbf{Startup is included.} Interpreter startup is \StartupShareMin--\StartupShareMax\% of each run. No region-of-interest markers exist yet.
  \item \textbf{One run per configuration.} Only the C smoke run has been repeated; sensitivity to small input or code-layout changes is untested.
  \item \textbf{One tiny workload.} A single integer loop says nothing about object-heavy or call-heavy code.
  \item \textbf{Hardware is a proposal.} The O3 baseline (decision D6) is not final, and no hardware sweep has been run on CPython.
  \item \textbf{Missing stdlib modules.} zlib, \texttt{\_decimal}, \texttt{\_ctypes}, \texttt{\_ssl} and others are not built, because their libraries are not cross-built for x86-64.
  \item \textbf{Builds are not release builds.} With \texttt{\BuildOpt} and no PGO or LTO, absolute speeds are below distributed CPython builds.
  \item \textbf{Simulation is slow.} O3 runs at about \SimMIPS M simulated instructions per host second, \HostMinMin--\HostMinMax{} minutes per run here.
\end{itemize}

\section{Next steps}

\begin{enumerate}
  \item Add region-of-interest markers (\texttt{m5\_work\_begin}/\texttt{m5\_work\_end}, already in the gem5 image) so startup is excluded, and move the JIT probe out of the measured region.
  \item Cross-build zlib, mpdecimal and libffi so standard benchmarks can import their modules.
  \item Choose a small benchmark set with fixed inputs and expected outputs that fits a per-run simulation budget (decision D5).
  \item Settle the baseline hardware and sweep values (D6), then run one-factor predictor and cache sweeps across all four engines.
\end{enumerate}

\appendix
\section{Runs used}

All numbers come from raw gem5 output and manifests under \texttt{results/raw/} (not committed), selected in \texttt{analysis/report\_runs.json}.
The CPython engine runs were made from repository state \RunsCommit.

\begin{table}[H]
\centering\footnotesize
\input{\GeneratedDir/tab-runs.tex}
\end{table}

\end{document}
